\section{Discussion}

Three observations stand out from the experimental evaluation. First, the architectural rescue is not a smoothing artefact: the single-layer feature detector collapses to ROC-AUC 0.7363 with a 26.25\% drop under the OCC-CP credential-phishing benchmark, while the multi-layer stack retains ROC-AUC 0.9572 with only a 4.12\% drop---a $\Delta$ of 22.13\% AUC points that comfortably satisfies the $\le 5\%$ robustness bound proposed in \S3.4. The ablation study in \S6.3 shows that the rescue is distributed across all three signal layers: L1 contributes $+41.0\%$ F1, L2 contributes $+13.9\%$ F1, L3 contributes $+22.5\%$ F1, and L3.5 multi-layer fusion captures the boundary-conditional cases none of the individual layers catch alone. The dominance of L1 is consistent with the intuition that OOD adversaries hide in the profiled identity surface, not in feature novelty \citep{gao2023dtgi, xiao2024sentinel, xiao2022robust}.

Second, matching insight \#2 from the Sailboat analogy in \S1, the \texttt{OOD\_CREDENTIAL\_PHISH} profile exposes a structural weakness in single-layer detectors that the multi-layer stack converts into a measurable architectural signal. Single-layer feature detectors operate on the assumption that benign traffic clusters tightly while insider traffic is observable as novelty---exactly the assumption the OOD adversary violates by operating during business hours with stolen credentials, low novelty scores, identity-proxy reuse, and removable-media transfers \citep{li2021threat, mavroeidis2018framework}. TURRET OS breaks this assumption by anchoring detection on cross-format file metadata (L1) rather than identity-conditioned novelty signal, gating the alert on joint rule-state through L3, and validating the joint evidence through L4. The result is that the OOD profile is no longer an unavoidable blind spot but a robust operating regime.

Third, the 100\% ISO/IEC 27043 attribute coverage demonstrated in \S6.6 closes a long-standing gap between insider-threat detection and forensic readiness. Prior deployed systems reported alerts only, leaving the inference-to-evidence chain to be reconstructed ad hoc during incident response \citep{aki2025provenance, shoderu2025dfrsaas, shoderu2026dfrconference}. TURRET OS commits every alert to a tamper-evident W3C-PROV record signed with Ed25519 over the Merkle root, providing Identification, Collection, Acquisition, Preservation, Analysis, Presentation, Chain of Custody, and Integrity Verification in a single deployable artefact \citep{yusof2025provcon}. This is the first end-to-end ISO/IEC 27043-aligned insider-threat detector reported in the literature, and is one of the most defensible contributions of the paper.

Practically, the deployment budget reported in \S6.7---15.0~s wall-clock per 1~K activity records, 9.4 MB on disk for the 35-user $\times$ 90-day D5 corpus, WORM-archived raw artefacts at 4.8 GB---makes TURRET OS deployable inside an existing SOC analyst workflow without specialised hardware. The 80-core, 503.58 GB RAM, NVIDIA L40S platform is well within mid-tier enterprise appliance capabilities, and the wall-clock budget is dominated by L3.5 GraphSAGE inference (6.3~s); therefore the marginal cost of adding the L1, L2, and L4 layers is small relative to the gains in robustness and admissibility.

The runtime safety invariant in \S5.6---that the top-5 SHAP feature importances must belong to a fixed operational whitelist---addresses a frequently overlooked leakage risk in graph-based insider detectors. On D5 the dominant feature is \texttt{copy\_to\_removable}, a semantically interpretable signal that confirms the cross-format metadata anchor is operational rather than artefact-driven \citep{li2021threat, mavroeidis2018framework}. Together with the chronological 60-30-10 split, this invariant makes TURRET OS difficult to game through feature-engineering attacks and audit-friendly during peer-review.

Finally, the comparison regime in \S6.2---ROC-AUC 0.9997, F1@0.5\%FPR 0.8774, MCC 0.9770 on D3---improves over the published baselines surveyed in \S2 (PS0 at ROC-AUC 0.92, F1 0.88; MEWRGNN at ROC-AUC 0.94, F1 0.87; DTGI at ROC-AUC 0.91, F1 0.86; TGCN-DA at ROC-AUC 0.95, F1 0.89; SENTINEL at ROC-AUC 0.93, F1 0.88) \citep{gao2023dtgi, li2023tgcnda, xiao2024sentinel, xiao2022robust}. Two distinctions matter: (i) the operating threshold is pinned to F1 at 0.5\% FPR---a precision-critical deployability constraint absent from the published baselines; (ii) the comparison is run on a single shared chronological 60-30-10 split, so the gains cannot be attributed to split-divergent tuning.
